\begin{proof}[Proof of Theorem \ref{LeMuFu}]
We argue in the beautiful framework exploited in Reutenauer's book \cite[1.5]{Reu}
using the complete tensor product
$$
\mathscr{A}=R\lala X,Y\rara \bar\otimes R\lala X,Y\rara
$$
equipped with a product induced from 
the shuffle product (resp. the concatenation product)
on the left (resp. right) of $\bar\otimes$. 
Recall that the ring of
$R$-linear endomorphisms $\mathrm{End}_RR\lala X,Y\rara$
can be embedded into $\mathscr{A}$ by 
$f\mapsto \sum_{w\in W} w\otimes f(w)$, and that the product of 
$\mathscr{A}$ restricts to the convolution product of $\mathrm{End}_RR\lala X,Y\rara$
defined by $f\ast g:=\mathrm{conc}\circ(f\otimes g)\circ \Delta$
(where ``$\mathrm{conc}$'' means concatenation of left and right sides of $\otimes$).
Note that, for $f\in\mathrm{End}_R R\lala X,Y\rara$ and  $J\in R\lala X,Y\rara$,
we have $f(J)=\sum_{w\in W} \lp w,J\rp f(w)$.

Since, by Corollary \ref{coro}, every word $w$ can be written as 
$\sum_{\bt\in \N_0^{(\infty)}}
\lp \SS^{(\bt)},w \rp \MM^{(\bt)}$, 
the element of $\mathscr{A}$ corresponding to 
the identity $\mathrm{id}\in \mathrm{End}_RR\lala X,Y\rara$ is:
\begin{align*}
\sum_{w\in W} w\otimes w &=
\sum_w w\otimes \sum_\bt \lp \SS^{(\bt)},w \rp\MM^{(\bt)} 
=
\sum_\bt (\sum_w  \lp\SS^{(\bt)},w\rp w) \otimes \MM^{(\bt)} \\
&= 
\sum_\bt \SS^{(\bt)} \otimes \MM^{(\bt)} \\
&=
\left(\sum_{d=0}^\infty 
\sum_{\btau\in\N_0^d} \SS^{(\btau;0)}\otimes \MM^{(\btau;0)}\right)
\cdot
\left(\sum_{t=0}^\infty X^t\otimes X^t\right),
\end{align*}
where we have used
$\SS^{(\bt)}=\SS^{(\btau;t)}=\SS^{(\btau;0)}\sha X^t$ and 
$\MM^{(\bt)}=\MM^{(\btau;t)}=\MM^{(\btau;0)}\cdot X^t$.
Observing that both factors of the last expression above correspond to specific 
$R$-linear endomorphisms,
we can apply $\mathrm{id}$ to $J$ as the convolution product of them and find from
$\Delta(J)=J\otimes J$ that
\begin{equation} \label{J-expansion}
J=\mathrm{id}(J)=
\left(\sum_{d=0}^\infty 
\sum_{\btau\in\N_0^d} \lp \SS^{(\btau;0)},J \rp \MM^{(\btau;0)}\right)
\left(\sum_{t=0}^\infty \frac{(c_X)^t}{t !} X^t\right).
\end{equation}
Note here that the pairing of $J$ with 
$X^t=X^{\sha\, t}/t!$ is equal to $(c_X)^t/t !$,
as easily seen from the fact that the specialization 
$J(X,0)\in R\lala X\rara$ at $Y=0$ is a group like
element $\exp(c_X\cdot X)$.
To complete the proof of Theorem \ref{LeMuFu},
given a fixed $\bk=(\bkappa;k_\infty)=(k_1,\dots,k_d; k_\infty)\in\N_0^{(\infty)}$
and $0\le s\le k_\infty$, 
we compute
the coefficient of~$w_{(\kappa;s)}=X^{k_1}Y\cdots X^{k_d}YX^s$ in the expansion 
of the first factor of the above right-hand side as follows:
\begin{align*}
&\sum_{d=0}^\infty 
\sum_{\btau\in\N_0^d} 
\lp \SS^{(\btau;0)},J \rp 
\lp \MM^{(\btau;0)},w_{(\bkappa;s)} \rp 
=\left\lp\sum_{d=0}^\infty 
\sum_{\btau\in\N_0^d} \lp\SS^{(\btau;0)},J \rp \MM^{(\btau;0)},
w_{(\bkappa;s)} \right\rp \\
=&\left\lp\sum_{d=0}^\infty 
\sum_{\btau\in\N_0^d} 
\biggl\lp\SS^{(\btau;0)},\sum_{\bu\in \N_0^{(\infty)}}(J,w_\bu)w_\bu\biggr\rp
 \MM^{(\btau;0)},
w_{(\bkappa;s)}\right\rp \\
=&\sum_{\bu} \lp J,w_\bu \rp \sum_{d=0}^\infty 
\sum_{\btau\in\N_0^d} \lp \SS^{(\btau;0)},w_\bu \rp
\lp \MM^{(\btau;0)}, w_{(\bkappa;s)} \rp .
\end{align*}
But since $\lp\SS^{(\btau;0)},w_\bu\rp \lp \MM^{(\btau;0)}, w_{(\bkappa;s)} \rp$ 
survives only when 
$\dep(\btau;0)=\dep(\bkappa;s)=\dep(\bu)$ and $|(\btau;0)|=|(\bkappa;s)|=|\bu|$,
the summation $\sum_\bu$ in the last expression above occurs only for 
those $\bu$ of the form 
$(\bkappa+\bs;0)\in \N_0^{(\infty)}$ 
with $\bs=(s_1,\dots,s_d)\in\Z^d$, 
$s=s_1+\cdots+s_d\ge 0$
(cf. also Remark \ref{rem2.6}).  
Then, it follows from Lemma \ref{lem2} that the last expression above is equal to 
$$
\sum_{d=0}^\infty \sum_{\substack{\bs\in\N_0^d \\ |(\bs;0)|=s}} 
\lp J, w_{(\bkappa+\bs;0)} \rp
(-1)^s
\binom{k_1+s_1}{k_1}\cdots \binom{k_d+s_d}{k_d}.
$$
(Note: the prescribed condition $\bs\in\Z^d$ has been replaced 
with $\bs\in\N_0^d$ because of the a posteriori survivals of binomial factors.)
From this and \eqref{J-expansion} together with
$\lp J, w_{(\bkappa+\bs;0)} \rp=c_{(k_1+s_1,\dots,k_d+s_d;0)}$, 
we conclude the assertion.
\end{proof}
